\scanpage{4}
i.e. un hom de schémas $\operatorname{Spec}K' - \operatorname{Spec}K$, on a un morphisme de Topos correspondant
\[
(1)\qquad B_{K'}-B_{K^{\unsure{\deleted{\prime}|*}}}
\]
associé à un hom de groupoïdes fond.
\[
(\unsure{2|1})\qquad \Pi_{K'}\longrightarrow\Pi_{K^{\unsure{\deleted{\prime}|*}}}.
\]
Ceci s'explicite en disant qu'\unsure{un|une} \deleted{\unsure{clôture}} objet de $\Pi_{K'}$ (i.e. $p^t$. de $B_{K'}$, ou rev. universel de $B_{K'}$, \deleted{)} ou clôture \unsure{sép.|\deleted{alg.}\inserted{sép.}}\inserted{$\bar K'$} de $K'$) en définit un de $\Pi_{K^{\unsure{\deleted{\prime}|*}}}$ (ainsi, on prend $\bar K^{\unsure{\deleted{\prime}|*}}=$ clôture \inserted{alg} \unsure{sép|\deleted{sép}|\deleted{alg}} de $K$ dans $\bar K'$) et pour deux $p^t$ correspondants, on a un hom de groupes fondamentaux correspondants, qui s'interpr\unsure{è|e}te p.ex. comme
\[
\unsure{(3)|3)|(n)}\qquad E_{\bar K'/K'}-E_{\bar K^{\unsure{\deleted{\prime}|*}}/K^{\unsure{\deleted{\prime}|*}}}
\]
et qui permet de réconstituer l'hom. de topos comme une ``restriction des scalaires''.

\margin{\unsure{1|$\vert$}}
L'image de (3) est le sous-groupe fermé de $E_{\bar K/K}$ qui correspond à la sous-\unsure{extension}
\margin{3}
